\documentclass[10pt,a4paper]{amsart}
\usepackage{amsmath,amssymb,amsthm,graphicx}
\usepackage{verbatim}
\newenvironment{reference}{\comment}{\endcomment}
%\newenvironment{reference}{}{}
\newenvironment{slogan}{\comment}{\endcomment}
\newenvironment{history}{\comment}{\endcomment}

% For commutative diagrams we use Xy-pic
\usepackage[all]{xy}

% We use 2cell for 2-commutative diagrams.
\xyoption{2cell}
% Theorem environments.
%
\theoremstyle{plain}
\newtheorem{theorem}[subsection]{Theorem}
\newtheorem{proposition}[subsection]{Proposition}
\newtheorem{lemma}[subsection]{Lemma}

\theoremstyle{definition}
\newtheorem{definition}[subsection]{Definition}
\newtheorem{example}[subsection]{Example}
\newtheorem{exercise}[subsection]{Exercise}
\newtheorem{situation}[subsection]{Situation}

\theoremstyle{remark}
\newtheorem{remark}[subsection]{Remark}
\newtheorem{remarks}[subsection]{Remarks}

\numberwithin{equation}{subsection}

% Macros
%
\def\lim{\mathop{\mathrm{lim}}\nolimits}
\def\colim{\mathop{\mathrm{colim}}\nolimits}
\def\Spec{\mathop{\mathrm{Spec}}}
\def\Hom{\mathop{\mathrm{Hom}}\nolimits}
\def\Ext{\mathop{\mathrm{Ext}}\nolimits}
\def\SheafHom{\mathop{\mathcal{H}\!\mathit{om}}\nolimits}
\def\SheafExt{\mathop{\mathcal{E}\!\mathit{xt}}\nolimits}
\def\Sch{\mathit{Sch}}
\def\Mor{\mathop{\mathrm{Mor}}\nolimits}
\def\Ob{\mathop{\mathrm{Ob}}\nolimits}
\def\Sh{\mathop{\mathit{Sh}}\nolimits}
\def\NL{\mathop{N\!L}\nolimits}
\def\CH{\mathop{\mathrm{CH}}\nolimits}
\def\proetale{{pro\text{-}\acute{e}tale}}
\def\etale{{\acute{e}tale}}
\def\QCoh{\mathit{QCoh}}
\def\Ker{\mathop{\mathrm{Ker}}}
\def\Im{\mathop{\mathrm{Im}}}
\def\Coker{\mathop{\mathrm{Coker}}}
\def\Coim{\mathop{\mathrm{Coim}}}

% Boxtimes
%
\DeclareMathSymbol{\boxtimes}{\mathbin}{AMSa}{"02}

%
% Macros for moduli stacks/spaces
%
\def\QCohstack{\mathcal{QC}\!\mathit{oh}}
\def\Cohstack{\mathcal{C}\!\mathit{oh}}
\def\Spacesstack{\mathcal{S}\!\mathit{paces}}
\def\Quotfunctor{\mathrm{Quot}}
\def\Hilbfunctor{\mathrm{Hilb}}
\def\Curvesstack{\mathcal{C}\!\mathit{urves}}
\def\Polarizedstack{\mathcal{P}\!\mathit{olarized}}
\def\Complexesstack{\mathcal{C}\!\mathit{omplexes}}
% \Pic is the operator that assigns to X its picard group, usage \Pic(X)
% \Picardstack_{X/B} denotes the Picard stack of X over B
% \Picardfunctor_{X/B} denotes the Picard functor of X over B
\def\Pic{\mathop{\mathrm{Pic}}\nolimits}
\def\Picardstack{\mathcal{P}\!\mathit{ic}}
\def\Picardfunctor{\mathrm{Pic}}
\def\Deformationcategory{\mathcal{D}\!\mathit{ef}}
\usepackage{cleveref}
\crefname{originalremark}{original remark}{original remarks}
\Crefname{originalremark}{Original remark}{Original remarks}
\makeatletter
\newlabel{remark-differential-object-selfmap}{{0AMK}{source}{}{sourcepointer-1}{}}
\newlabel{remark-differential-object-selfmap@cref}{{[originalremark][0][]0AMK}{[1][0][]source}{}{}{}}
\makeatother
\title{The spectral sequence associated to a filtered complex}
\author{The Stacks Project Authors}
\date{stacks-project:homology.tex:local-snapshot-20261012; GNU FDL 1.2 or later}
\begin{document}
\maketitle
\noindent\textit{Editorial note.} All delivered mathematical text is the author's own, copied from the pinned Stacks Project chapter file homology.tex; no statement or proof step was rewritten, reordered or removed.\par\medskip
\noindent\textit{Editorial note.} The proof is the author's complete argument. It opens with the author's own reading recommendation 'we strongly suggest the reader skip this proof', which by LANE\_BRIEF 5 is a suggestion to the reader and not a refusal to prove; the general construction is then carried out in full inside the counted span and the closing sentences derive the formulae for the boundaries and cocycles.\par\medskip
\noindent\textit{Editorial note.} Every labelled reference occurring inside the counted spans is carried in assembly.external\_references as a non-load-bearing navigation pointer (same-chapter ones at their definition site, cross-chapter ones at the invocation inside the counted span), so no load-bearing mathematics is replaced by a stub. No bibliography citation occurs inside the counted spans.\par\medskip
\noindent\textit{Editorial note.} The same-name check for this card scanned every staged Stacks chapter of this repository -- all sources/stacks-*/pinned-original-source.tex files, the card's own chapter excluded -- and the labels found are recorded in dedup.same\_name\_cross\_chapter on the card.\par\medskip
\noindent\textit{Editorial note.} The preamble is the chapter's own native preamble excerpt already used for this source: the comment-type reference, slogan and history environments (native preamble lines 8-13), the Xy-pic setup (native preamble lines 14-18) and the native theorem declarations with the mathematical macros (native preamble lines 154-217, the actual last line of the file).\par\medskip
\noindent\textit{Allowed standard prerequisites (editorial).}\begin{itemize}
\item Spectral sequences from an injective self-map of a differential object. Same chapter, Remark remark-differential-object-selfmap. Supplies the construction of a spectral sequence from an injective morphism of differential objects, which is the device that produces the spectral sequence in the proof, together with the descriptions of the boundaries and cocycles used at the end.
\item Filtered complexes and the associated graded. Same chapter, the definitions of a filtered complex, of the filtration steps F\textasciicircum{}p and of the graded pieces gr\textasciicircum{}p. Supplies the objects and the differential whose graded pieces are identified with the first two pages.
\item The countable direct sum case of the construction. Same chapter, the discussion immediately preceding the lemma. The proof refers to it in its opening sentence for the case where the abelian category has exact countable direct sums, and then gives the general construction.
\end{itemize}\medskip
\ifdefined\hypertarget\hypertarget{sourcepointer-1}{}\fi
\noindent\textit{Source reference.} Original remark 0AMK: sources/stacks-homology/pinned-original-source.tex, source lines 5811--5811. Navigation pointer invoked at proof line 6295; the counted span contains the pointer only and the cited result is not reproduced inside it, so the pointer is registered as a non-load-bearing navigation stub in the established form for Stacks material. No bibliography citation is involved.\par\medskip
\begin{lemma}
\label{lemma-spectral-sequence-filtered-complex}
Let $\mathcal{A}$ be an abelian category. Let $(K^\bullet, F)$ be a
filtered complex of $\mathcal{A}$. There is a spectral sequence
$(E_r, d_r)_{r \geq 0}$ in the category of bigraded objects of $\mathcal{A}$
associated to $(K^\bullet, F)$ such that $d_r$ has bidegree $(r, - r + 1)$
and such that $E_r$ has bigraded pieces $E_r^{p, q}$ and maps
$d_r^{p, q} : E_r^{p, q} \to E_r^{p + r, q - r + 1}$ as given above.
Furthermore, we have $E_0^{p, q} = \text{gr}^p(K^{p + q})$,
$d_0^{p, q} = \text{gr}^p(d^{p + q})$,
and $E_1^{p, q} = H^{p + q}(\text{gr}^p(K^\bullet))$.
\end{lemma}

\begin{proof}
If $\mathcal{A}$ has countable direct sums and if countable direct
sums are exact, then this follows from the discussion above.
In general, we proceed as follows; we strongly suggest the reader
skip this proof. Consider the bigraded object $A = (F^{p + 1}K^{p + 1 + q})$
of $\mathcal{A}$, i.e., we put $F^{p + 1}K^{p + 1 + q}$ in degree $(p, q)$
(the funny shift in numbering to get numbering correct later on).
We endow it with a differential $d : A \to A[0, 1]$ by using $d$
on each component. Then $(A, d)$ is a differential bigraded object.
Consider the map
$$
\alpha : A \to A[-1, 1]
$$
which is given in degree $(p, q)$ by the inclusion
$F^{p + 1}K^{p + 1 + q} \to F^pK^{p + 1 + q}$.
This is an injective morphism of differential objects
$\alpha : (A, d) \to (A, d)[-1, 1]$. Hence, we can apply
Remark \ref{remark-differential-object-selfmap}
with $S = [0, 1]$ and $T = [1, -1]$.
The corresponding spectral sequence $(E_r, d_r)_{r \geq 0}$
of bigraded objects is the spectral sequence we are looking
for. Let us unwind the definitions a bit. First of all we have
$E_r = (E_r^{p, q})$. Then, since $T^rS = [r, -r + 1]$
we have $d_r : E_r \to E_r[r, -r + 1]$ which means that
$d_r^{p, q} : E_r^{p, q} \to E_r^{p + r, q - r + 1}$.

\medskip\noindent
To see that the description of the graded pieces hold, we argue
as above. Namely, first we have
$$
E_0 = \Coker(\alpha : A \to A[-1, 1])[0, -1] =
\Coker(\alpha[0, -1] : A[0, -1] \to A[-1, 0])
$$
and by our choice of numbering above this gives
$$
E_0^{p, q} = \Coker(F^{p + 1}K^{p + q} \to F^pK^{p + q}) = \text{gr}^pK^{p + q}
$$
The first differential is given by
$d_0^{p, q} = \text{gr}^pd^{p + q} : E_0^{p, q} \to E_0^{p, q + 1}$.
Next, the description of the boundaries $B_r$ and the cocycles $Z_r$
in Remark \ref{remark-differential-object-selfmap}
translates into a straightforward manner into the formulae
for $Z_r^{p, q}$ and $B_r^{p, q}$ given above.
\end{proof}

\end{document}
